

Linguistic innovation is often used by subgroups in society. 

“We were the people who were not in the papers. We lived in the blank white spaces at the edges of print. It gave us more freedom.
We lived in the gaps between the stories.”

“It's impossible to say a thing exactly the way it was, because of what you say can never be exact, you always have to leave something out, there are too many parts, sides, crosscurrents, nuances; too many gestures, which could mean this or that, too many shapes which can never be fully described, too many flavors, in the air or on the tongue, half-colors, too many.”

Swearing in certain contexts is fine and in others, it's norm breaking, as are nudity, sex, and bodily fluids. Is swearing a norm violation because it is obscene, like a couple lustily copulating in a McDonalds booth or a pile of shit on the floor of an elevator, or because it is ungrammatical?

Imagine swearing isn't about the shock value of certain words, but instead exists as a secret language hidden within our own. It has its own twisted grammar, a code of expletives and crude turns of phrase that flout everything your schoolteacher drilled into you. Just like slang, or the cryptic lingo of a trade, swearing becomes a badge of belonging.

Those within the circle have mastered this rebellious grammar. They hurl a curse with casual precision, knowing exactly when it adds emphasis or earns knowing laughter.  But those on the outside fumble, either too shocked to use the words at all or dropping them clumsily into a sentence, revealing their outsider status. It's not about morality, but about fluency in this forbidden dialect.

This hidden language enforces its own sense of order. A well-timed profanity among familiar faces sparks camaraderie – its violation of “polite” rules becomes a wink of shared understanding. Yet, that same word unleashed against a stranger is an attack, precisely because they exist outside the imagined tribe where this twisted grammar is allowed. Swearing doesn't signal a failure of vocabulary, but rather a pointed choice to speak – or not speak – the subversive language that draws lines between those in the know, and those firmly excluded.